\section{Introduccion}

El trabajo practico estudia sistemas mediante simulacion computacional y comparacion con resultados teoricos cuando estos existen. El primer modelo analizado es una cola M/M/1, implementada mediante una simulacion de eventos discretos en Python. En una etapa posterior se agregara el modelo de inventario y la comparacion con AnyLogic.

Para el modelo de colas se busca observar como cambia la congestion al variar la relacion entre tasa de arribo y tasa de servicio. Ademas, se analizan colas con capacidad finita para medir la probabilidad de denegacion de servicio cuando el sistema no puede aceptar nuevos clientes.

Los parametros principales se leen desde \texttt{config/default\_parameters.json}, lo que permite modificar la configuracion experimental sin cambiar el codigo fuente.
